\section{Related Work}

\label{sec:related_work}

Memory control and organisation for language-model agents has attracted several architectural proposals. MemGPT introduces an operating-system-inspired virtual context manager that moves information across memory tiers, demonstrating document analysis and multi-session chat scenarios \citep{ref_959d6858d034c6e1}. The present study's fixed L1 Rail does not implement MemGPT's autonomous virtual-memory controller; MemGPT serves only as conceptual context for the tiered-storage vocabulary. A-MEM creates structured notes and dynamically links and evolves memories following a Zettelkasten-inspired organisation \citep{ref_1aacba116cbe3c5b}. Our fixed lexical ranking lacks autonomous linking and evolution, and no empirical comparison to A-MEM was run. Generative Agents synthesise natural-language experience records into reflections and dynamically retrieve them for planning behaviour in an interactive simulation \citep{ref_6e953134a3c04402}. This study neither implements reflection or planning nor tests social behaviour; it draws only on the general notion of retrieving stored records at query time.

Evidence access and position effects constitute a second theme. Lost in the Middle demonstrates that relevant information position affects multi-document question answering and key-value retrieval in long contexts, with middle positions often harder to exploit \citep{ref_9dbf664b6954efbc}. That finding motivates attention to how information is placed within a prompt, but it does not establish that the prefix truncation observed in our baseline is caused by the same attentional phenomenon; the baseline failure is deterministic removal of text beyond the budget, not degraded attention to retained text. MemoryBank supports retrieval and updates for sustained personalised interaction, modelling forgetting and reinforcement influenced by elapsed time and importance \citep{ref_713e886950406893}. Our temporal-decay term is a simpler heuristic and does not replicate MemoryBank's companion evaluation or its reinforcement dynamics.

Evaluation of long-term interactive memory forms the third cluster. LongMemEval benchmarks chat assistants on extraction, multi-session reasoning, temporal reasoning, knowledge updates, and abstention, explicitly separating indexing, retrieval, and reading stages \citep{ref_ab2c226739f2f52a}. The submitted adaptation covers only a narrow single-session-user extraction slice converted into fixed fact episodes; it does not exercise multi-session reasoning, temporal reasoning, updates, or abstention, and the resulting scores are not official LongMemEval benchmark results. Together, these six sources frame the design space within which the present two-arm comparison sits, while none is replicated or empirically benchmarked here.
